%
% Combining high-quality radio Xmax reconstructions and muon reconstructions for mass composition (chapter 6)
%
\chapter{Combined Analysis of Muon and Electromagnetic Reconstructions}\label{chap:radiomuon}
Description AERA Shower maximum reconstructions on the SD energy scale (Bjarni work) and the reconstruction of the absolute EM energy from AERA

The UMD provides direct muon measurements of air showers on a per-event basis, which in principle can be combined with per-event radio $X_\text{max}$ reconstructions to increase the overall mass sensitivity of Auger, both per-event and for the average mass composition.
The AERA and UMD footprints are co-located making simultaneous measurements feasible; however, implementation is difficult due to both geometrical and timeline restrictions for each detector system.
The geometry of UMD scintillator panels is optimal for vertical air-shower reconstructions while the spacing of the AERA radio antennas makes the AERA $X_{\rm max}$ reconstruction more efficient for air showers with $\theta_{\rm zen}\,>\,$\qty{40}{\degree} (see~\Cref{fig:aera_umd_histograms}).
Regarding timeline constraints, AERA, recording air showers since 2011, has fulfilled its primary science goal of proving the usefulness of air-shower reconstructions in the \qtyrange{30}{80}{\mega\hertz} band near Malarg\"{u}e and may eventually be decomissioned.
The UMD did not obtain a stable reconstruction until January 2018, with deployment only finished in March 2026, severely limiting the data overlap between the two detectors.
The radio $X_{\rm max}$ reconstruction will eventually be ported to the RD antennas, deployed at each Auger SD, yet the spacing of RD will aim to reconstruct even more highly-inclined air showers ($\theta_{\text{zen}}>$\,\qty{65}{\degree}).

\begin{figure}
  \centering
  \includegraphics[trim={0cm 0cm 0cm 0cm},clip,width=0.49\textwidth]{chapter6/AERA_BiasFreeXmax_ZenithHistogram.pdf}
  \includegraphics[trim={0cm 0cm 0cm 0cm},clip,width=0.49\textwidth]{chapter6/MDSD750_ICRC2025_test7_reconstructions_ZenithHistogram.pdf}
  \caption[Reconstructed zenith angle histograms for AERA and UMD.]{Reconstructed zenith angle histograms from the bias free sample of the AERA $X_\text{max}$ reconstructions (left) and the standard UMD + SD-750 reconstruction of the Pierre Auger Observatory (right). The label ``icrc2025-test7'' refers to the version tag for the Auger observatory's reconstruction software. Both reconstructions shown here are for Auger Phase I data only.}
  \label{fig:aera_umd_histograms}
\end{figure}

\section{Muon and Shower Maximum Reconstructions with an Electromagnetic Energy}\label{sec:radioenergyscale}
Comparison of ln(A) Plots (avg. and sigma)

\section{Combination of Muon and Shower Maximum Observables into a Single Observable}\label{sec:combo}
Is this possible?

\section{Discussion and Interpretation}\label{sec:interpretation}
